\documentclass{amsart}
\usepackage{amsmath,amssymb,amsthm,hyperref}
\theoremstyle{plain}
\newtheorem{theorem}{Theorem}
\newtheorem{lemma}{Lemma}
\newtheorem{proposition}{Proposition}
\newtheorem{corollary}{Corollary}
\theoremstyle{definition}
\newtheorem{remark}{Remark}
\begin{document}

\title{On the existence of a second Hamiltonian cycle in $4$-regular Hamiltonian graphs}
\maketitle

\section{Statement and status of the problem}

We are asked to prove or disprove the following.

\begin{quote}
\emph{For every finite simple $4$-regular graph $G$ that contains a Hamiltonian
cycle, and for every Hamiltonian cycle $C$ of $G$, there is a Hamiltonian cycle
$C'\neq C$ of $G$; equivalently, every $4$-regular Hamiltonian graph has at
least two distinct Hamiltonian cycles.}
\end{quote}

This statement is precisely \emph{Sheehan's Conjecture} (J.~Sheehan, 1975),
which asserts that every Hamiltonian $4$-regular simple graph contains a second
Hamiltonian cycle.  Despite intensive work it remains \textbf{open} in general:
no proof and no counterexample is known.  Consequently we cannot present a
complete unconditional proof of the displayed statement, nor can we disprove it.
We are honest about this throughout.

What we \emph{can} do, and do rigorously and completely below, is:
\begin{enumerate}
\item[(i)] Prove the underlying parity theorem of Thomason from first principles
(the ``lollipop'' argument), Theorem~\ref{thm:thomason}.  This is the tool that
settles the odd-degree analogue of the problem completely.
\item[(ii)] Deduce from it a clean, broad \emph{sufficient condition}
(Theorem~\ref{thm:main}) under which a $4$-regular Hamiltonian graph provably has
a second Hamiltonian cycle.  In particular this proves the conjecture for every
$4$-regular Hamiltonian graph $G$ that has a Hamiltonian cycle $C$ whose
complementary $2$-factor $G-E(C)$ is bipartite (equivalently, has no odd cycle):
Corollary~\ref{cor:bipartite}.
\item[(iii)] Explain precisely why the parity method, which resolves all
odd-degree cases, fails for the even-regular case, thereby locating exactly the
obstruction that makes the full conjecture hard
(Remark~\ref{rem:obstruction}).
\end{enumerate}

Throughout, ``graph'' means finite, and Hamiltonian cycles are counted as
\emph{edge sets} (unoriented, no distinguished starting point); two Hamiltonian
cycles are the same iff they have the same edge set.

\section{Thomason's lollipop theorem}

\begin{theorem}[Thomason]\label{thm:thomason}
Let $H$ be a finite graph and let $e=xy$ be an edge of $H$ such that every vertex
of $H$ other than $x$ and $y$ has odd degree.  Then the number of Hamiltonian
cycles of $H$ that contain $e$ is even.
\end{theorem}

\begin{proof}
We may assume $H$ has at least three vertices (otherwise there are no
Hamiltonian cycles and the statement holds vacuously).

\medskip
\noindent\textbf{Configurations.}
Call a sequence $P=(x,y,u_3,u_4,\dots,u_m)$ a \emph{configuration} if it is a
Hamiltonian path of $H$ (so $m=|V(H)|$, the $u_i$ are pairwise distinct, they
exhaust $V(H)$, consecutive vertices are adjacent in $H$) whose first two
vertices are $x,y$ in this order.  Thus a configuration is a Hamiltonian path
with the fixed first edge $e=xy$ traversed from $x$ to $y$; its \emph{free end}
is its last vertex $w:=u_m$, and its \emph{root edge} is $e$.  Let $\mathcal P$
be the (finite) set of all configurations.

A configuration $P$ with free end $w$ is called \emph{closing} if $w$ is adjacent
to $x$ in $H$ (note $w\neq x$ since $x$ is the first vertex, and $w\neq y$ unless
$m=2$, which we excluded).  We record:

\medskip
\noindent\textbf{Claim 1: Closing configurations biject with Hamiltonian cycles
through $e$.}
Indeed, if $C$ is a Hamiltonian cycle containing $e=xy$, delete $e$ from $C$ to
obtain a Hamiltonian path with endpoints $x$ and $y$; orienting it from $x$
gives a unique sequence $(x,y,\dots,w)$ where $w$ is the neighbour of $x$ on $C$
other than $y$.  This is a configuration, and it is closing because $w\sim x$.
Conversely, a closing configuration $(x,y,\dots,w)$ yields the Hamiltonian cycle
obtained by adding the edge $wx$; this contains $e$.  The two maps are mutually
inverse: starting from a Hamiltonian cycle through $e$, the construction returns
the same cycle, and starting from a closing configuration we recover it.  (Each
Hamiltonian cycle through $e$ has exactly two edges at $x$, namely $e$ and one
other; only the traversal starting $x\to y$ produces a configuration whose root
edge is $e$, so the correspondence is one-to-one.)  This proves Claim~1.

\medskip
\noindent\textbf{The lollipop graph $L$.}
We define a graph $L$ on vertex set $\mathcal P$.  Let $P=(x,y,\dots,t,w)\in
\mathcal P$ have free end $w$ and let $t$ be the vertex preceding $w$ in $P$.
For each neighbour $u$ of $w$ in $H$ with
\[
u\neq t \quad\text{and}\quad u\neq x,
\]
the vertex $u$ occurs somewhere in $P$, say $P=(x,y,\dots,u,s,\dots,w)$ where $s$
is the successor of $u$ in $P$ (such $s$ exists because $u\neq w$).  Define the
\emph{rotation of $P$ at $u$} to be the sequence
\[
\rho_u(P)=(x,y,\dots,u,\,w,\dots,s),
\]
obtained by reversing the segment of $P$ strictly after $u$.  Concretely,
$\rho_u(P)$ deletes the edge $us$ and adds the edge $uw$; it is again a
Hamiltonian path with the same first two vertices $x,y$ (the prefix up to and
including $u$ is unchanged, in particular the root edge $e$ is preserved because
$u\neq x$), hence $\rho_u(P)\in\mathcal P$, and its new free end is $s$.

Put an (undirected) edge of $L$ between $P$ and $\rho_u(P)$ for each admissible
$u$.  This is well defined as an undirected graph: if $Q=\rho_u(P)$ has free end
$s$ and predecessor (the vertex before $s$ in $Q$) equal to $w$, then $u$ is a
neighbour of $s$ in $Q$ with $u\neq w$ and $u\neq x$, and applying the rotation
of $Q$ at $u$ deletes $uw$ and adds $us$, recovering $P$.  Thus the relation
``$Q=\rho_u(P)$ for some admissible $u$'' is symmetric, and $L$ is a genuine
(simple) graph; distinct admissible $u$ give distinct neighbours $\rho_u(P)$
because they have distinct free ends $s$ (the successor of $u$ along $P$ is
determined by $u$).

\medskip
\noindent\textbf{Claim 2: The degree of $P$ in $L$.}
With $w,t$ as above,
\[
\deg_L(P)=\#\{\,u\in N_H(w): u\neq t,\ u\neq x\,\}
=\deg_H(w)-1-[\,w\sim x\,],
\]
where $[\,\cdot\,]$ is the Iverson bracket: we subtract $1$ for the predecessor
$t$ (which is always a neighbour of $w$ and is never equal to $x$ because, for
$m\ge 3$, $t$ is the vertex before the last and $x$ is the first), and we
subtract one more exactly when $x$ itself is a neighbour of $w$ (then $u=x$ is
excluded as it is not allowed to use the root vertex for a rotation).  This is
Claim~2.

\medskip
\noindent\textbf{Parity of degrees.}
Recall the hypothesis: every vertex of $H$ other than $x,y$ has odd degree.

\emph{If $P$ is not closing}, then $w\not\sim x$, so by Claim~2
$\deg_L(P)=\deg_H(w)-1$.  Here the free end $w$ satisfies $w\neq x$ and $w\neq y$
(if $w=y$ then $m=2$, excluded).  Hence $w$ is one of the ``other'' vertices, so
$\deg_H(w)$ is odd, and $\deg_L(P)=\deg_H(w)-1$ is \textbf{even}.

\emph{If $P$ is closing}, then $w\sim x$, so $\deg_L(P)=\deg_H(w)-2$.  Again
$w\notin\{x,y\}$ (a closing configuration has $w\sim x$, and $w=y$ is impossible
for $m\ge3$), so $\deg_H(w)$ is odd and $\deg_L(P)=\deg_H(w)-2$ is \textbf{odd}.

Therefore the vertices of $L$ of \emph{odd} degree are \emph{exactly} the closing
configurations.

\medskip
\noindent\textbf{Conclusion.}
In any finite graph the number of vertices of odd degree is even (handshake
lemma).  Applied to $L$, this says the number of closing configurations is even.
By Claim~1 the number of closing configurations equals the number of Hamiltonian
cycles of $H$ that contain $e$.  Hence that number is even.
\end{proof}

\begin{corollary}[Thomason's second–Hamiltonian–cycle theorem]\label{cor:odd}
Let $H$ be a Hamiltonian graph in which every vertex has odd degree.  Then $H$
has at least two distinct Hamiltonian cycles.  In particular every cubic
($3$-regular) Hamiltonian graph has at least two Hamiltonian cycles.
\end{corollary}

\begin{proof}
Let $C$ be a Hamiltonian cycle of $H$ and let $e\in C$ be any edge.  Apply
Theorem~\ref{thm:thomason} with this $e=xy$: the hypothesis ``all vertices other
than $x,y$ have odd degree'' holds because \emph{all} vertices of $H$ have odd
degree.  Thus the number of Hamiltonian cycles through $e$ is even.  It is at
least $1$ because $C$ is one of them; being even, it is at least $2$.  Hence
there is a Hamiltonian cycle $C'\neq C$ containing $e$.  Cubic graphs have all
degrees equal to $3$, which is odd, so the last sentence follows.
\end{proof}

\section{A sufficient condition for $4$-regular graphs}

We now use Theorem~\ref{thm:thomason} to handle a large class of $4$-regular
Hamiltonian graphs.  The point is that, although a $4$-regular graph has all
degrees even (so Theorem~\ref{thm:thomason} does not apply to $G$ itself), we may
sometimes \emph{delete a matching} from the complementary $2$-factor to reach a
graph whose degrees are all odd except at the two endpoints of one cycle edge,
to which Theorem~\ref{thm:thomason} then applies.

\begin{lemma}\label{lem:2factor}
Let $G$ be a $4$-regular graph and $C$ a Hamiltonian cycle of $G$.  Then
$F:=G-E(C)$ (the graph on $V(G)$ whose edges are the edges of $G$ not in $C$) is
$2$-regular; hence $F$ is a vertex-disjoint union of cycles spanning $V(G)$.
\end{lemma}

\begin{proof}
Every vertex $v$ has degree $4$ in $G$ and degree exactly $2$ in $C$ (a cycle
uses two edges at each vertex), so $v$ has degree $4-2=2$ in $F$.  A finite graph
in which every vertex has degree $2$ is a vertex-disjoint union of cycles
covering all vertices.
\end{proof}

\begin{theorem}\label{thm:main}
Let $G$ be a $4$-regular graph and $C$ a Hamiltonian cycle of $G$, with
complementary $2$-factor $F=G-E(C)$.  Suppose there exist two vertices $x,y$ with
$xy\in E(C)$ and a matching $M\subseteq E(F)$ that saturates every vertex of $G$
except $x$ and $y$ (i.e.\ the set of vertices not covered by $M$ is exactly
$\{x,y\}$).  Then $G$ has a Hamiltonian cycle $C'\neq C$.
\end{theorem}

\begin{proof}
Define $H:=G-M$, the graph obtained from $G$ by deleting the edges of $M$.
We verify the hypotheses of Theorem~\ref{thm:thomason} for $H$ and the edge
$e=xy$.

\emph{$H$ is a simple graph and $e\in E(H)$.}  Deleting edges from a simple graph
keeps it simple.  Since $M\subseteq E(F)$ and $E(F)$ is disjoint from $E(C)$, the
edge $e=xy\in E(C)$ is not in $M$, hence $e\in E(H)$.

\emph{$C$ is a Hamiltonian cycle of $H$.}  All edges of $C$ lie in $E(G)\setminus
E(F)\supseteq E(G)\setminus M$, so none of them is deleted; thus
$E(C)\subseteq E(H)$ and $C$ is a Hamiltonian cycle of $H$.

\emph{Degree parities in $H$.}  Each vertex $v$ has $\deg_G(v)=4$.  Deleting $M$
lowers $\deg(v)$ by the number of edges of $M$ at $v$.  Since $M$ is a matching,
this number is $0$ or $1$.  By hypothesis $M$ covers every vertex except $x,y$;
so for $v\notin\{x,y\}$ exactly one edge of $M$ is at $v$, giving
$\deg_H(v)=4-1=3$, which is odd; and for $v\in\{x,y\}$ no edge of $M$ is at $v$,
giving $\deg_H(v)=4$.  Hence every vertex of $H$ other than $x,y$ has odd degree.

By Theorem~\ref{thm:thomason} applied to $H$ and $e=xy$, the number of
Hamiltonian cycles of $H$ containing $e$ is even.  The cycle $C$ is one such
(it is a Hamiltonian cycle of $H$ and $e\in E(C)$), so this even number is at
least $2$; hence there is a Hamiltonian cycle $C'$ of $H$ with $e\in E(C')$ and
$C'\neq C$.

Finally, $H$ is a subgraph of $G$ on the same vertex set, so every Hamiltonian
cycle of $H$ is also a Hamiltonian cycle of $G$.  Thus $C'$ is a Hamiltonian
cycle of $G$ different from $C$.
\end{proof}

\begin{corollary}\label{cor:bipartite}
Let $G$ be a $4$-regular graph having a Hamiltonian cycle $C$ such that the
complementary $2$-factor $F=G-E(C)$ is bipartite (equivalently, contains no odd
cycle).  Then $G$ has at least two distinct Hamiltonian cycles.
\end{corollary}

\begin{proof}
By Lemma~\ref{lem:2factor}, $F$ is a disjoint union of cycles.  Bipartite means
all these cycles have even length.  An even cycle $C_{2k}$ has a perfect matching
(take alternate edges); taking the union of one perfect matching from each cycle
of $F$ yields a perfect matching $M_0\subseteq E(F)$ of $G$, i.e.\ $M_0$ saturates
\emph{every} vertex.

Set $H:=G-M_0$.  Then, exactly as in the proof of Theorem~\ref{thm:main}, every
edge of $C$ survives in $H$ (since $M_0\subseteq E(F)$ is disjoint from $E(C)$),
so $C$ is a Hamiltonian cycle of $H$; and each vertex $v$ now has
$\deg_H(v)=4-1=3$, so $H$ is a cubic graph.  Cubic graphs have all degrees odd,
so by Corollary~\ref{cor:odd}, $H$ has at least two Hamiltonian cycles.  Since
$H\subseteq G$ on the same vertex set, $G$ has at least two distinct Hamiltonian
cycles.

(Alternatively, $M_0$ saturates everyone; pick any edge $e=xy\in E(C)$ and let
$M=M_0\setminus\{$the two $M_0$-edges at $x$ and at $y\}$ to reduce to
Theorem~\ref{thm:main}; but the direct cubic argument above is cleaner.)
\end{proof}

\begin{remark}
Theorem~\ref{thm:main} is strictly more general than
Corollary~\ref{cor:bipartite}: it also covers some graphs whose complementary
$2$-factor $F$ has odd cycles, namely whenever $F$ admits a matching whose two
uncovered vertices are joined by an edge of $C$.  For instance this happens if
$F$ has exactly two odd cycles, one through $x$ and one through $y$ for some edge
$xy\in C$, and all other cycles of $F$ are even.
\end{remark}

\section{Where the method stops, and the open core}

\begin{remark}[The obstruction to a full proof]\label{rem:obstruction}
The parity argument of Theorem~\ref{thm:thomason} is the standard route to
``second Hamiltonian cycle'' results, and it settles \emph{every} graph in which
all relevant vertices have odd degree (Corollary~\ref{cor:odd}).  For a genuinely
$4$-regular graph $G$ \emph{all} degrees are even ($=4$), so the theorem does not
apply to $G$ directly.  The repair used above is to delete a matching to create
odd degrees while keeping $C$ intact; this requires the complementary
$2$-factor $F=G-E(C)$ to contain a matching missing only the two endpoints of
some edge of $C$.  Such a matching need not exist: if $F$ has an odd cycle, that
cycle cannot be perfectly matched, and the parities of the cycles of $F$ may
forbid leaving uncovered exactly a pair of $C$-adjacent vertices.  Concretely,
one cannot always reduce a $4$-regular Hamiltonian graph to a cubic (or
all-odd-degree) Hamiltonian graph by deleting a matching while preserving the
given cycle $C$.

More fundamentally, run the lollipop construction of
Theorem~\ref{thm:thomason} directly on $G$ (which has $\deg_G(w)=4$ for the free
end $w$ of any configuration).  By Claim~2 of that proof, a closing
configuration then has $L$-degree $\deg_G(w)-2=2$ and a non-closing one has
$L$-degree $\deg_G(w)-1=3$.  Thus the \emph{odd}-degree vertices of the lollipop
graph $L$ are the \emph{non}-closing configurations, and the handshake lemma
gives no information about the number of closing configurations (= Hamiltonian
cycles through $e$).  Indeed one checks on examples that for $4$-regular graphs
the number of Hamiltonian cycles through a fixed edge can be odd, so no
edge-parity statement of Thomason type can hold.  This is exactly why the
even-regular case escapes the parity method.
\end{remark}

\begin{remark}[Status of the full statement]
The displayed problem is Sheehan's Conjecture and is open.  It is known to hold
for $r$-regular Hamiltonian graphs once $r$ is large enough (Thomassen; with the
best current bounds due to subsequent work of Haxell--Seamone--Verstra\"ete and
Girão--Kittipassorn--Narayanan), and trivially for all \emph{odd} $r$ by
Corollary~\ref{cor:odd}; but for the specific value $r=4$ no complete proof is
known.  Accordingly, the results above
(Theorems~\ref{thm:thomason},~\ref{thm:main} and
Corollaries~\ref{cor:odd},~\ref{cor:bipartite}) constitute the strongest
unconditional statements we can establish: a complete proof of the underlying
parity theorem, and a proof of the conjecture for the broad family of
$4$-regular Hamiltonian graphs possessing a Hamiltonian cycle whose
complementary $2$-factor is bipartite (and slightly beyond).  We make no claim
to have resolved the conjecture in full, and we have exhibited no counterexample
(none is known); hence the problem can be neither proved nor disproved here.
\end{remark}

\end{document}
